\section{Введение}\label{sec:intro}

Мы рассматриваем задачу минимизации суммы функций
\[
 f(x)=\frac1n\sum_{i=1}^n f_i(x)\to\min_{x\in\R^d}.
\]
Задачи такого вида возникают при минимизации эмпирического риска
в машинном обучении, где $f_i$ --- функция потерь на $i$-м объекте
обучающей выборки \citep{shalev2014understanding,bottou2018optimization}.
В федеративном и распределённом обучении слагаемое $f_i$ --- это функция
потерь на данных $i$-го устройства или вычислительного узла
\citep{mcmahan2017communication}. Градиент одного слагаемого часто вычисляется
существенно дешевле полного градиента. Поэтому сложность методов мы измеряем
числом обращений к оракулу, который возвращает значение и градиент одного слагаемого.

Для суммы $L$-гладких функций число обращений к оракулу позволяют сократить
методы уменьшения дисперсии \citep{johnson2013svrg,xiao2014proxsvrg}.
Они периодически вычисляют полный градиент $f$, а на остальных
шагах используют градиенты отдельных слагаемых $f_i$.
Нижние оценки для произвольных рандомизированных методов показывают,
какие зависимости от числа слагаемых и точности улучшить нельзя
\citep{woodworth2016finite}.

Однако условие $L$-гладкости не выполняется, например, для экспоненциальных
функций потерь. Условие $(L_0,L_1)$-гладкости введено в работе
\citet{zhang2019gradient}, чтобы объяснить успех методов
с клиппированием градиента. При обучении языковой модели на основе LSTM авторы
заметили, что локальная константа гладкости растёт вместе с нормой градиента.
Аналогичная зависимость обнаружена и при обучении трансформеров
\citep{crawshaw2022robustness}.
Условие $(L_0,L_1)$-гладкости как раз допускает такой рост кривизны.
Его обобщения и методы первого порядка для таких функций изучались
в работах
\citet{li2023generalized,gorbunov2024methods,vankov2025optimizing,tyurin2025unified}.
Для конечных сумм нужно учитывать одновременно и число слагаемых, и рост кривизны.
Например, в пуассоновской регрессии с логарифмической функцией связи
\citep{stanPoisson}
$f_i(x)=\exp(a_i^{\mathsf T}x)-b_i a_i^{\mathsf T}x$, где $b_i\ge0$,
гессиан не ограничен на всём пространстве, но
\[
 \norm{\nabla^2 f_i(x)}
 \le b_i\norm{a_i}^2+\norm{a_i}\norm{\nabla f_i(x)}.
\]
В задаче восстановления фазы по измерениям интенсивности используются потери
$f_i(x)=(\ip{a_i}x^2-b_i)^2$
\citep{candes2015wirtinger}. Здесь также приходится минимизировать конечную сумму
слагаемых с неограниченной кривизной.
Ещё один пример --- экспоненциальные потери в задаче классификации
\citep[пример~1.5]{gorbunov2024methods}.
Во всех этих задачах при большом числе наблюдений вычислять полный градиент
может быть дорого, а сам градиент не является глобально липшицевым.
Поэтому нужны методы, которые учитывают рост кривизны и не вычисляют
градиенты всех слагаемых на каждом шаге. Ниже предполагается,
что функция $f$ достигает минимума, а расстояние до точки минимума, начальная невязка
и неоднородность слагаемых ограничены.

\paragraph{Особенность суммы функций.}
В рассматриваемой задаче каждое слагаемое удовлетворяет неравенству
\[
 0\preceq\nabla^2f_i(x)
 \preceq\bigl(L_{0,i}+L_1\norm{\nabla f_i(x)}\bigr)I.
\]
Даже в точке минимума среднего градиенты отдельных слагаемых могут быть
большими: равенство $\nabla f(x^\ast)=0$ означает лишь, что их сумма
равна нулю. Поэтому малость нормы полного градиента не влечёт малости
констант гладкости слагаемых, от которых зависят оценки методов уменьшения
дисперсии. Какова оракульная сложность минимизации суммы в этой ситуации?
Мы отвечаем на этот вопрос в настоящей работе.

\paragraph{Обозначения.} Через $\varepsilon>0$ обозначается требуемая точность решения задачи по функции: метод должен выдать точку $\hat x$, для которой $f(\hat x)-f^\ast\le\varepsilon$. Знак $\norm{\cdot}$ всюду означает евклидову норму вектора и спектральную норму матрицы, $[x,y]$ --- отрезок с концами $x,y\in\R^d$, $B(x,r)$ --- замкнутый шар радиуса $r$ с центром $x$. Сумму первых $k$ членов последовательности мы обозначаем соответствующей заглавной буквой: например, $A_k=a_1+\dots+a_k$ и $A_0=0$.


Для рандомизированных методов указанная точность требуется
в математическом ожидании. Все проекции евклидовы; через $\Pi_C$
обозначается проекция на непустое замкнутое выпуклое множество $C$.

\paragraph{Основной результат.}
Пусть слагаемые выпуклы, дважды непрерывно дифференцируемы и ограничены
снизу, а $n^{-1}\sum_iL_{0,i}\le L_0$. Начальную точку $x_0$ будем считать
нулевой: этого можно добиться сдвигом, относительно которого класс
задач инвариантен. Предполагается, что $f$ достигает минимума в точке
$x^\ast$, причём $\norm{x^\ast}\le R$ и $f(0)-f^\ast\le\Delta$.
При условиях теоремы~\ref{thm:main} верхняя и нижняя оценки сложности
имеют, с точностью до логарифмических множителей, порядок
\[
 n\bigl(1+(L_1R)^{2/3}\bigr)
 +\sqrt{\frac{n(L_0+L_1^2\delta)R^2}{\varepsilon}+n(L_1R)^2}.
\]
Здесь под сложностью понимается математическое ожидание числа запросов
значений и градиентов отдельных слагаемых в худшем случае,
достаточного для выполнения $\E[f(\hat x)-f^\ast]\le\varepsilon$.
Нижняя оценка справедлива для произвольных рандомизированных методов;
размерность задачи при этом не ограничена.
Точное определение класса и логарифмические множители приведены
в разделе~\ref{sec:setting}.

\paragraph{Замечание о логарифмических множителях.}
В верхней оценке скрыт множитель $\log^2N$, где $N$
из~\eqref{eq:complexity-parameters} растёт линейно по $\Delta/\varepsilon$.
В худшем случае, когда $\Delta$ экспоненциально велико по $L_1R$,
этот логарифм может иметь порядок $L_1R$, так что множитель $\log^2N$ может иметь порядок $(L_1R)^2$,
то есть давать дополнительный полиномиальный множитель.

Слагаемое $\sqrt{nL_0R^2/\varepsilon}$ известно из теории минимизации
суммы $L$-гладких функций \citep{woodworth2016finite}.
Слагаемое $n(L_1R)^{2/3}$ отражает число обращений к оракулу,
необходимое, чтобы приблизиться к минимуму из области, где градиенты
экспоненциально велики. Слагаемое
\[
 \sqrt{\frac{nL_1^2R^2\delta}{\varepsilon}+n(L_1R)^2}
\]
возникает из-за того, что кривизна слагаемых вблизи минимума
определяется их градиентами, которые могут оставаться большими даже в точке
минимума среднего.

\paragraph{Метод и эксперименты.}
Верхнюю оценку достигает метод YETI. При малой локальной кривизне
он выполняет ускоренный шаг, решая вспомогательную задачу методом SVRG,
а при большой --- проецирует вспомогательную точку на полупространство,
содержащее точку минимума. Кроме того, YETI использует эвристику
с ограниченным числом запросов к оракулу: она позволяет воспользоваться
удачными предложенными точками, не ухудшая теоретическую оценку.
В экспериментах YETI сравнивается с другими методами и проверяется
его критерий остановки.

\paragraph{Связь с известными методами.}
YETI использует ускоренную проксимальную схему
Монтейро--Свайтера \citep{ms2013} и идею адаптивного выбора
регуляризации без решения неявного уравнения на каждом шаге
\citep{carmon2022adaptive}. Локальная задача решается
проекционным вариантом SVRG. Ускорение за счёт приближённого решения
регуляризованных задач изучалось также в схеме Catalyst
\citep{lin2015catalyst,lin2018catalyst} и в работе
\citet{frostig2015}. В YETI выбор шага зависит от вычисляемой оценки
кривизны суммы: если эта оценка слишком велика, метод вместо шага строит
отделяющее полупространство. Оценка сложности учитывает запросы
и локального решателя, и эвристики.

Нижняя оценка опирается на принцип Яо \citep{yao1977}: трудная задача
выбирается случайно из конечного семейства, после чего проводится усреднение
по случайным величинам алгоритма. Скрытые координаты и распределение
информации по разным слагаемым продолжают известный подход к построению
нижних оценок для задач минимизации суммы \citep{woodworth2016finite}.
Для рассматриваемого класса мы применяем экспоненциальные
преобразования выпуклых цепочек и выбираем их параметры так,
чтобы выполнялись ограничения на $L_1$, $\delta$ и начальную невязку.

Настоящая статья продолжает нашу предыдущую работу \citet{rudomanov2026everest}, в которой
оракульная сложность этого класса найдена для одной функции ($n=1$):
её порядок равен $\sqrt{L_0R^2/\varepsilon}+L_1R$, а при малой
начальной невязке --- $\sqrt{L_0R^2/\varepsilon}+(L_1R)^{2/3}$.
Здесь рассматривается сумма функций, а оракул возвращает значение
и градиент лишь одного слагаемого, поэтому для получения верхней оценки
мы используем уменьшение дисперсии.
Результаты для одной функции
\citep{vankov2025optimizing,tyurin2026accelerated} относятся
к оракулу, вычисляющему полный градиент. В частности, оптимальность
метода \citet{tyurin2026accelerated} при достаточно малых $\varepsilon$ основана
на классической нижней оценке для гладкого подкласса и касается только
ведущего члена $\sqrt{L_0R^2/\varepsilon}$. Оптимальность слагаемого
$n(L_1R)^{2/3}$ для оракула, возвращающего градиент одного слагаемого,
из неё не вытекает.
Аналогично,
оценки для стохастических методов в случае общей точки минимума всех слагаемых
\citep{gorbunov2024methods}
не охватывают произвольную положительную неоднородность $\delta$.

\paragraph{Организация работы.}
Постановка задачи и основной результат приведены в разделе~\ref{sec:setting},
метод YETI и его анализ --- в разделах~\ref{sec:cheap}--\ref{sec:cheap-analysis},
основные эксперименты --- в разделе~\ref{sec:experiments}
(дополнительные эксперименты и их детали --- в приложении~\ref{app:experiments}), нижняя оценка ---
в разделе~\ref{sec:lower}. В основном тексте приведены
результаты и схема доказательства нижней оценки, а вспомогательные
доказательства и варианты метода вынесены в приложения. Сводная таблица
известных оценок приведена в приложении~\ref{app:literature}.
